\section{abstract}
    Phishing continues to grow because it targets the individual rather than the system. Furthermore, blacklists are inherently reactive—by the time a new site is reported, it has already achieved its objective. To address this, the project proposes an AI agent that evaluates the phishing risk of a given URL without relying on such lists. Instead, it analyzes characteristics of the URL and its domain—such as length, the presence of unusual characters, the use of an IP address rather than a domain name, and the domain's registration age. When a site is accessible, the agent also examines content elements, including login forms or visual similarities to well-known brands. The model employs XGBoost, Random Forest, and LightGBM to assess the relative importance of various features, retaining only the most significant ones. This refined subset of features serves as the foundation for a two-level stacking architecture: XGBoost and Random Forest each generate a phishing probability, which then feeds into a simpler, second-level XGBoost model responsible for the final prediction. The system is trained on a dataset of real phishing and legitimate URLs gathered from public sources, allowing the agent to process raw URLs directly rather than relying on pre-calculated features. The goal is to provide a functional agent, while also offering clarity regarding its limitations specifically, how well it generalizes to new, previously unseen sites and how susceptible it might be to evasion tactics using specially crafted URLs.
\newline \par

\textbf{Keywords:} Phishing, Cybersecurity, Artificial Intelligence, Machine Learning, URL Analysis, Risk Assessment, Ensemble Learning, Explainability.
